% =============================================================================
% Spaces, contours and spectral-sequence pages: constructors
% =============================================================================
% Building blocks for topology and complex-analysis pictures. A constructor
% takes the data of an object and draws it; the objects themselves are files
% under figures/objects/: spaces/ (spheres, balls, cylinder, torus, its
% fundamental square, coffee cup, Moebius band), contours/ (integration
% contours) and diagrams/ (commutative diagrams, a tikz-cd diagram in a node).
% \input by dzg-constructors.tex, so "@" is a letter here.
%
% One decision per former pic or command:
%   sphere=<n>, ball=<n>, cylinder, torus, torus fundamental square, coffee
%     cup, \mobiusband: one particular space each -> objects/spaces/.
%   quarter disc contour, strip rectangle contour: one particular contour
%     each -> objects/contours/.
%   \kleinquotients, \enriquescorrespondence, \cowedgesquare,
%     \monaddescentaxioms: one particular diagram per variant -> objects/diagrams/.
%   space labels, \spacelabel: constructor, the canonical label of any space.
%   \complexaxes: constructor, the axes of C from their extents.
%   ss page, \ssentry, \ssdifferential: constructor, a page from its ranges.
%
% Commands:
%   \spacelabel{<x,y>}{<label>}  the node -label at (x,y), if space labels
%   \complexaxes{<xmin>}{<xmax>}{<ymin>}{<ymax>}
%                                the coordinate axes of C, and the label C at
%                                the top right corner if space labels
%   \pic (E) [<ss keys>] {ss page}, \ssentry, \ssdifferential   see below
%
% Keys:
%   space labels=canonical|none  canonical labels of spaces and contours
%
% Conditional for objects: \ifspacelabels.

\newif\ifspacelabels \spacelabelstrue
\tikzset{
  space labels/.is choice,
  space labels/canonical/.code={\spacelabelstrue},
  space labels/none/.code={\spacelabelsfalse},
}
\newcommand{\spacelabel}[2]{\ifspacelabels\node (-label) at (#1) {#2};\fi}
\newcommand{\complexaxes}[4]{%
  \draw[coordinate axis] (#1,0) -- (#2,0);
  \draw[coordinate axis] (0,#3) -- (0,#4);
  \ifspacelabels\node[font=\small] at (#2-0.3,#4-0.3) {$\mathbb{C}$};\fi}

% -----------------------------------------------------------------------------
% A page of a spectral sequence. The cell (p,q) is the coordinate <name>-p-q at
% (p * ss column sep, q * ss row sep); cells exist for p, q up to three beyond
% the page, so differentials can leave it. The coordinate axes run between
% p = -1 and p = 0 and between q = -1 and q = 0; a page that starts at 0 gets an
% L-shaped frame. Keys (on the pic):
%   ss p={0,...,5}, ss q={0,...,5}   the columns and rows of the page
%   ss p labels={0,n}, ss q labels=  tick labels in place of the indices
%   ss p name=$p$, ss q name=$q$     axis names; ss ticks=false hides ticks
%   ss column sep=2cm, ss row sep=1cm
%   ss homological                   d^r has bidegree (-r, r-1); the default is
%                                    cohomological, d_r of bidegree (r, 1-r)
% Figures place entries and differentials:
%   \ssentry{E}{p}{q}{$E_2^{p,q}$}   the node E-p-q, replacing the coordinate
%   \ssdifferential[<draw options>]{E}{r}{p}{q}
%                                    d_r from (p,q), labelled d_r^{p,q} (or
%                                    d^r_{p,q}); `ss unlabelled` omits the label
% -----------------------------------------------------------------------------
\def\dzg@ssP{0,...,5}\def\dzg@ssQ{0,...,5}
\def\dzg@ssPlabels{}\def\dzg@ssQlabels{}
\def\dzg@ssPname{$p$}\def\dzg@ssQname{$q$}
\def\dzg@ssconv{coh}
\newif\ifdzg@ssticks \dzg@sstickstrue
\newif\ifdzg@sslabel \dzg@sslabeltrue
\tikzset{
  ss p/.store in=\dzg@ssP, ss q/.store in=\dzg@ssQ,
  ss p labels/.store in=\dzg@ssPlabels, ss q labels/.store in=\dzg@ssQlabels,
  ss p name/.store in=\dzg@ssPname, ss q name/.store in=\dzg@ssQname,
  ss column sep/.initial=2cm, ss row sep/.initial=1cm,
  ss ticks/.is if=dzg@ssticks,
  ss homological/.code={\def\dzg@ssconv{hom}},
  ss unlabelled/.code={\dzg@sslabelfalse},
}
% The first and last entries of a foreach list.
\def\dzg@ssends#1#2#3{%
  \def#2{}\foreach \dzg@i in #1 {\ifx#2\empty\xdef#2{\dzg@i}\fi\xdef#3{\dzg@i}}}

\tikzset{pics/ss page/.style={/tikz/transform shape, code={\tikzset{transform shape=false}%
  \pgfmathsetlengthmacro\dzg@sx{\pgfkeysvalueof{/tikz/ss column sep}}%
  \pgfmathsetlengthmacro\dzg@sy{\pgfkeysvalueof{/tikz/ss row sep}}%
  \global\let\dzg@ssx\dzg@sx \global\let\dzg@ssy\dzg@sy
  \dzg@ssends{\dzg@ssP}\dzg@pmin\dzg@pmax
  \dzg@ssends{\dzg@ssQ}\dzg@qmin\dzg@qmax
  \pgfmathtruncatemacro\dzg@pa{\dzg@pmin-3}\pgfmathtruncatemacro\dzg@pb{\dzg@pmax+3}%
  \pgfmathtruncatemacro\dzg@qa{\dzg@qmin-3}\pgfmathtruncatemacro\dzg@qb{\dzg@qmax+3}%
  \foreach \p in {\dzg@pa,...,\dzg@pb} \foreach \q in {\dzg@qa,...,\dzg@qb}
    \coordinate (-\p-\q) at (\p*\dzg@sx, \q*\dzg@sy);
  \expandafter\xdef\csname dzg@ssconv@\pgfkeysvalueof{/tikz/name prefix}\endcsname{\dzg@ssconv}%
  % axes
  \pgfmathsetmacro\dzg@xl{min(\dzg@pmin,0)-0.5}\pgfmathsetmacro\dzg@xr{\dzg@pmax+0.6}%
  \pgfmathsetmacro\dzg@yb{min(\dzg@qmin,0)-0.5}\pgfmathsetmacro\dzg@yt{\dzg@qmax+0.6}%
  \draw (\dzg@xl*\dzg@sx, -0.5*\dzg@sy) -- (\dzg@xr*\dzg@sx, -0.5*\dzg@sy)
    node[right] {\dzg@ssPname};
  \draw (-0.5*\dzg@sx, \dzg@yb*\dzg@sy) -- (-0.5*\dzg@sx, \dzg@yt*\dzg@sy)
    node[above] {\dzg@ssQname};
  \ifdzg@ssticks
    \foreach \p [count=\k] in \dzg@ssP {%
      \ifx\dzg@ssPlabels\empty\def\dzg@t{$\p$}\else
        \foreach \l [count=\j] in \dzg@ssPlabels {\ifnum\j=\k\xdef\dzg@t{$\l$}\fi}\fi
      \node[font=\small, below] at (\p*\dzg@sx, \dzg@yb*\dzg@sy) {\dzg@t};}
    \foreach \q [count=\k] in \dzg@ssQ {%
      \ifx\dzg@ssQlabels\empty\def\dzg@t{$\q$}\else
        \foreach \l [count=\j] in \dzg@ssQlabels {\ifnum\j=\k\xdef\dzg@t{$\l$}\fi}\fi
      \node[font=\small, left] at (\dzg@xl*\dzg@sx, \q*\dzg@sy) {\dzg@t};}
  \fi}}}

\newcommand{\ssentry}[4]{\node (#1-#2-#3) at (#1-#2-#3) {#4};}
\newcommand{\ssdifferential}[5][]{%
  \begingroup
  \dzg@sslabeltrue
  \tikzset{#1}%
  \edef\dzg@c{\csname dzg@ssconv@#2\endcsname}\def\dzg@hom{hom}%
  \ifx\dzg@c\dzg@hom
    \pgfmathtruncatemacro\dzg@tp{#4-#3}\pgfmathtruncatemacro\dzg@tq{#5+#3-1}%
    \def\dzg@lab{$d^{#3}_{#4,#5}$}%
  \else
    \pgfmathtruncatemacro\dzg@tp{#4+#3}\pgfmathtruncatemacro\dzg@tq{#5-#3+1}%
    \def\dzg@lab{$d_{#3}^{#4,#5}$}%
  \fi
  \ifdzg@sslabel
    \draw[map arrow, shorten <=1pt, shorten >=1pt, #1] (#2-#4-#5) -- (#2-\dzg@tp-\dzg@tq)
      node[midway, fill=white, inner sep=1pt, font=\small] {\dzg@lab};
  \else
    \draw[map arrow, shorten <=1pt, shorten >=1pt, #1] (#2-#4-#5) -- (#2-\dzg@tp-\dzg@tq);
  \fi
  \endgroup}
